% Options for packages loaded elsewhere
\PassOptionsToPackage{unicode}{hyperref}
\PassOptionsToPackage{hyphens}{url}
\documentclass[
  ignorenonframetext,
]{beamer}
\newif\ifbibliography
\usepackage{pgfpages}
\setbeamertemplate{caption}[numbered]
\setbeamertemplate{caption label separator}{: }
\setbeamercolor{caption name}{fg=normal text.fg}
\beamertemplatenavigationsymbolsempty
% remove section numbering
\setbeamertemplate{part page}{
  \centering
  \begin{beamercolorbox}[sep=16pt,center]{part title}
    \usebeamerfont{part title}\insertpart\par
  \end{beamercolorbox}
}
\setbeamertemplate{section page}{
  \centering
  \begin{beamercolorbox}[sep=12pt,center]{section title}
    \usebeamerfont{section title}\insertsection\par
  \end{beamercolorbox}
}
\setbeamertemplate{subsection page}{
  \centering
  \begin{beamercolorbox}[sep=8pt,center]{subsection title}
    \usebeamerfont{subsection title}\insertsubsection\par
  \end{beamercolorbox}
}
% Prevent slide breaks in the middle of a paragraph
\widowpenalties 1 10000
\raggedbottom
\AtBeginPart{
  \frame{\partpage}
}
\AtBeginSection{
  \ifbibliography
  \else
    \frame{\sectionpage}
  \fi
}
\AtBeginSubsection{
  \frame{\subsectionpage}
}
\usepackage{iftex}
\ifPDFTeX
  \usepackage[T1]{fontenc}
  \usepackage[utf8]{inputenc}
  \usepackage{textcomp} % provide euro and other symbols
\else % if luatex or xetex
  \usepackage{unicode-math} % this also loads fontspec
  \defaultfontfeatures{Scale=MatchLowercase}
  \defaultfontfeatures[\rmfamily]{Ligatures=TeX,Scale=1}
\fi
\usepackage{lmodern}
\ifPDFTeX\else
  % xetex/luatex font selection
\fi
% Use upquote if available, for straight quotes in verbatim environments
\IfFileExists{upquote.sty}{\usepackage{upquote}}{}
\IfFileExists{microtype.sty}{% use microtype if available
  \usepackage[]{microtype}
  \UseMicrotypeSet[protrusion]{basicmath} % disable protrusion for tt fonts
}{}
\makeatletter
\@ifundefined{KOMAClassName}{% if non-KOMA class
  \IfFileExists{parskip.sty}{%
    \usepackage{parskip}
  }{% else
    \setlength{\parindent}{0pt}
    \setlength{\parskip}{6pt plus 2pt minus 1pt}}
}{% if KOMA class
  \KOMAoptions{parskip=half}}
\makeatother
\usepackage{longtable,booktabs,array}
\usepackage{caption}
\captionsetup[table]{skip=6pt}
\newcounter{none} % for unnumbered tables
\usepackage{calc} % for calculating minipage widths
\usepackage{caption}
% Make caption package work with longtable
\makeatletter
\def\fnum@table{\tablename~\thetable}
\makeatother
\setlength{\emergencystretch}{3em} % prevent overfull lines
\providecommand{\tightlist}{%
  \setlength{\itemsep}{0pt}\setlength{\parskip}{0pt}}
% Manuscript macro/package fallbacks (newcommand -> providecommand).
% Core mathematical packages
\usepackage{amsmath,amssymb,amsthm}
\usepackage{mathtools}
% Permit bibliography and theorem prose to break without overfull boxes.
% Algorithm formatting
% Tables
\usepackage{booktabs}
\usepackage{multirow}
% Graphics
% Cross-referencing with smart naming
% Configure cleveref for custom environments
% Theorem environments with consistent numbering
% Code listings used by supplementary implementation examples
\usepackage{listings}
\lstset{basicstyle=\ttfamily\small,breaklines=true,columns=fullflexible}
% End of theorem environment declarations
% Math operators
\DeclareMathOperator*{\argmax}{arg\,max}
\DeclareMathOperator*{\argmin}{arg\,min}
\DeclareMathOperator{\sign}{sign}
\DeclareMathOperator{\KL}{KL}
\DeclareMathOperator{\Tr}{Tr}
% Custom commands for notation consistency
\providecommand{\calA}{\mathcal{A}}
\providecommand{\calB}{\mathcal{B}}
\providecommand{\calC}{\mathcal{C}}
\providecommand{\calD}{\mathcal{D}}
\providecommand{\calF}{\mathcal{F}}
\providecommand{\calG}{\mathcal{G}}
\providecommand{\calH}{\mathcal{H}}
\providecommand{\calI}{\mathcal{I}}
\providecommand{\calM}{\mathcal{M}}
\providecommand{\calO}{\mathcal{O}}
\providecommand{\calP}{\mathcal{P}}
\providecommand{\calS}{\mathcal{S}}
\providecommand{\calT}{\mathcal{T}}
\providecommand{\calW}{\mathcal{W}}
% Note: \Phi is a standard LaTeX command, do not redefine
% Trust notation
\providecommand{\trust}[2]{\mathcal{T}_{#1 \to #2}}
\providecommand{\trustt}[3]{\mathcal{T}_{#1 \to #2}^{#3}}
% Belief notation
\providecommand{\belief}[2]{\mathcal{B}_{#1}(#2)}
\providecommand{\belieft}[3]{\mathcal{B}_{#1}^{#2}(#3)}
% Cognitive state
\providecommand{\cogstate}[1]{\sigma_{#1}}
% Attack notation
\providecommand{\attack}[1]{\mathcal{A}_{#1}}
\providecommand{\adversary}[1]{\Omega_{#1}}
% Defense notation
\providecommand{\firewall}{\mathcal{F}}
\providecommand{\sandbox}{\mathcal{S}_{box}}
% Probability and expectation
\providecommand{\E}{\mathbb{E}}
\providecommand{\Prob}{\mathbb{P}}
\providecommand{\indicator}{\mathbb{1}}
% QED symbol
\providecommand{\qedsymbol}{$\blacksquare$}
% Page Layout (Slightly smaller margins)
% Hyperlink Styling
% List of Figures and List of Tables in TOC

% Slide layout defaults for warning-clean scientific decks.
\usepackage{etoolbox}
\IfFileExists{xurl.sty}{\usepackage{xurl}}{}
\IfFileExists{seqsplit.sty}{\usepackage{seqsplit}}{\newcommand{\seqsplit}[1]{#1}}
\protected\def\breakseq#1{\seqsplit{#1}}
\protected\def\breaktt#1{\begingroup\ttfamily\seqsplit{#1}\endgroup}
\setlength{\emergencystretch}{6em}
\tolerance=5000
\hbadness=10000
\hfuzz=1pt
\setlength{\tabcolsep}{2pt}
\AtBeginEnvironment{longtable}{\tiny\renewcommand{\arraystretch}{0.86}\setlength{\tabcolsep}{1pt}}
\AtBeginEnvironment{tabular}{\tiny\renewcommand{\arraystretch}{0.86}\setlength{\tabcolsep}{1pt}}
\AtBeginEnvironment{equation}{\tiny}
\AtBeginEnvironment{equation*}{\tiny}
\AtBeginEnvironment{align}{\tiny}
\AtBeginEnvironment{align*}{\tiny}
\AtBeginEnvironment{itemize}{\footnotesize}
\AtBeginEnvironment{enumerate}{\footnotesize}
\AtBeginEnvironment{description}{\footnotesize}
\setbeamerfont{caption}{size=\tiny}
\setbeamerfont{caption name}{size=\tiny}
\setbeamerfont{normal text}{size=\small}
\setbeamerfont{frametitle}{size=\small}
\setbeamerfont{section title}{size=\footnotesize}
\setbeamerfont{subsection title}{size=\footnotesize}
\setbeamertemplate{section page}{%
  \centering
  \begin{beamercolorbox}[sep=12pt,center,wd=\paperwidth]{section title}
    \parbox{0.86\paperwidth}{\centering\usebeamerfont{section title}\insertsection\par}
  \end{beamercolorbox}
}
\setbeamertemplate{subsection page}{%
  \centering
  \begin{beamercolorbox}[sep=8pt,center,wd=\paperwidth]{subsection title}
    \parbox{0.86\paperwidth}{\centering\usebeamerfont{subsection title}\insertsubsection\par}
  \end{beamercolorbox}
}
\setlength{\abovecaptionskip}{2pt}
\setlength{\belowcaptionskip}{0pt}

% Natbib and cross-reference fallbacks — slides don't load natbib
% or cleveref, but combined-PDF manuscript prose may emit these
% commands. The fallback renders citations as a bracketed key list
% and cross-references as detokenized labels so slides don't fail on
% undefined control sequences or raw underscores. \providecommand is
% a no-op if packages load later.
\providecommand{\citep}[1]{[#1]}
\providecommand{\citet}[1]{#1}
\providecommand{\citealp}[1]{#1}
\providecommand{\citeauthor}[1]{#1}
\providecommand{\citeyear}[1]{#1}
\providecommand{\cref}[1]{\texttt{\detokenize{#1}}}
\providecommand{\Cref}[1]{\texttt{\detokenize{#1}}}
\providecommand{\crefrange}[2]{\texttt{\detokenize{#1}}--\texttt{\detokenize{#2}}}
\providecommand{\Crefrange}[2]{\texttt{\detokenize{#1}}--\texttt{\detokenize{#2}}}
\providecommand{\paragraph}[1]{\textbf{#1}\ }

% Additional theorem-like environments from manuscript preamble.
\newtheorem{observation}{Observation}
\newtheorem{formalization}{Formalization}
\newtheorem{principle}{Principle}

\newtheorem{proposition}{Proposition}
\newtheorem{hypothesis}{Hypothesis}
\newtheorem{remark}[theorem]{Remark}
\newtheorem{axiom}{Axiom}
\newtheorem{property}{Property}
\usepackage{bookmark}
\IfFileExists{xurl.sty}{\usepackage{xurl}}{} % add URL line breaks if available
\urlstyle{same}
% fallback for those not using the hyperref driver hyperxmp:
\makeatletter
\@ifundefined{xmpquote}{\newcommand{\xmpquote}[1]{#1}}{}
\makeatother
\hypersetup{
  hidelinks,
  pdfcreator={LaTeX via pandoc}}

\author{\texorpdfstring{}{}}
\date{}

\begin{document}

\begin{frame}
\newpage
\end{frame}

\section{Supplementary: Eusocial Insect Intelligence and Colony
Cognitive Security}\label{sec:eusocial-cogsec}

\subsection{Overview}\label{sec:eusocial-overview}

\begin{frame}[allowframebreaks]{Overview}
This supplementary material introduces \emph{colony cognitive security}
as a complementary paradigm to single-agent AI safety and alignment.
While the main CIF framework (\cref{sec:formal-framework}) addresses
cognitive integrity at the individual agent level, eusocial insect
colonies---ants, bees, termites---demonstrate that security properties
can emerge from collective dynamics that are irreducible to individual
behavior. This section formalizes these collective phenomena, identifies
the benchmark gap for multiagent cognitive security, and proposes
evaluation scenarios grounded in biological precedent.
\end{frame}

\subsubsection{The Paradigm Gap}\label{sec:paradigm-gap}

\begin{frame}[allowframebreaks]{The Paradigm Gap}
Contemporary AI security research exhibits a pronounced single-agent
bias. Existing benchmarks---jailbreak resistance, prompt injection
detection, harmful content refusal---evaluate individual models in
isolation ({Perez et al.} 2023; Wei et al. 2023). Even recent multiagent
security work (\cref{sec:threat-model}) often frames attacks as
adversary-versus-agent rather than adversary-versus-colony.

This mirrors a historical bias in behavioral biology. For decades,
researchers explained insect societies as aggregations of individual
optimizers, missing the insight that colonies function as
\emph{superorganisms} with collective cognition that transcends
individual capacity (Wilson 1971). The colony's cognitive
architecture---its ability to solve problems, allocate resources, and
respond to threats---emerges from interaction patterns, not individual
intelligence.

\end{frame}

\begin{frame}[allowframebreaks]{The Paradigm Gap}
\begin{observation}[Colony vs. Agent Security]
\label{obs:colony-agent}
Let $\mathcal{O} = \langle \mathcal{A}, \mathcal{C}, \mathcal{S}, \mathcal{P}, \Gamma \rangle$ be a multiagent operator. *Agent-level security* ensures that for all $a_i \in \mathcal{A}$, the cognitive state $\sigma_i$ remains within acceptable bounds. *Colony-level security* ensures that the collective function $\mathcal{F}_{\text{colony}}(\{\sigma_i\}_{i=1}^n)$ remains within acceptable bounds—even when individual $\sigma_i$ may be compromised.
\end{observation}

\end{frame}

\begin{frame}[allowframebreaks]{The Paradigm Gap}
These are orthogonal concerns. A colony can exhibit collective
resilience despite individual failures (Byzantine fault tolerance), and
conversely, individually secure agents can produce collectively
pathological outcomes (emergent misalignment).
\end{frame}

\begin{frame}
\end{frame}

\subsection{Theoretical Foundations}\label{sec:eusocial-theory}

\subsubsection{Stigmergy: Environment-Mediated
Coordination}\label{sec:stigmergy}

\begin{frame}[allowframebreaks]{Stigmergy: Environment-Mediated
Coordination}
Eusocial insects coordinate through \emph{stigmergy}---indirect
communication via environmental modification (Grassé 1959). Ants deposit
pheromones; bees perform waggle dances; termites build structures that
guide subsequent building. The environment becomes an external memory
and communication channel.

\begin{definition}[Stigmergic Operator]
\label{def:stigmergic-operator}
A *stigmergic operator* extends $\mathcal{O}$ with an environmental state $\mathcal{E}$:
\begin{equation}
\label{eq:stigmergic-operator}
\mathcal{O}_\Sigma = \langle \mathcal{A}, \mathcal{C}, \mathcal{S}, \mathcal{P}, \Gamma, \mathcal{E}, \Sigma \rangle
\end{equation}
where $\mathcal{E}(t): \mathcal{L} \times \mathcal{M} \to \mathbb{R}^+$ maps locations $l \in \mathcal{L}$ and marker types $m \in \mathcal{M}$ to signal intensities, and $\Sigma: \mathcal{A} \times \mathcal{E} \to \mathcal{E}'$ is the stigmergic update function.
\end{definition}

\end{frame}

\begin{frame}[allowframebreaks]{Stigmergy: Environment-Mediated
Coordination}
In AI systems, stigmergic analogs include:

\end{frame}

\begin{frame}[allowframebreaks]{Stigmergy: Environment-Mediated
Coordination}
\begin{itemize}
\tightlist
\item
  \textbf{Shared memory/state} --- Redis caches, vector databases, file
  systems
\item
  \textbf{Message queues} --- Kafka topics, RabbitMQ exchanges
\item
  \textbf{Artifact trails} --- Git commits, audit logs, provenance
  chains
\item
  \textbf{Embedding spaces} --- Semantic markers in shared vector stores
\end{itemize}

\end{frame}

\begin{frame}[allowframebreaks]{Stigmergy: Environment-Mediated
Coordination}
The critical insight is that attacks on \(\mathcal{E}\) constitute
attacks on the colony's cognitive substrate---analogous to the
\emph{cyberphysical niche} where AI agents operate.

\begin{definition}[Cyberphysical Niche]
\label{def:cyberphysical-niche}
The *cyberphysical niche* $\mathcal{N}$ of a stigmergic operator is the tuple:
\begin{equation}
\label{eq:niche}
\mathcal{N} = \langle \mathcal{E}, \mathcal{I}_{\text{ext}}, \mathcal{R}, \mathcal{H}_{\text{env}} \rangle
\end{equation}
where $\mathcal{I}_{\text{ext}}$ is the external information environment (web, APIs, sensors), $\mathcal{R}$ is the resource landscape (compute, memory, tokens), and $\mathcal{H}_{\text{env}}$ is environmental history.
\end{definition}
\end{frame}

\subsubsection{Emergent Collective Function}\label{sec:emergence}

\begin{frame}[allowframebreaks]{Emergent Collective Function}
Colony-level computation arises from simple individual rules applied in
parallel. Ant foraging, bee nest-site selection, and termite mound
construction all exhibit problem-solving capacity that exceeds any
individual's cognitive capacity (Bonabeau et al. 1999).

\begin{definition}[Emergent Collective Function]
\label{def:emergent-function}
For a stigmergic operator $\mathcal{O}_\Sigma$ with agents $\mathcal{A} = \{a_1, \ldots, a_n\}$, the *emergent collective function* $\mathcal{F}_c$ is:
\begin{equation}
\label{eq:emergent-function}
\mathcal{F}_c: \mathcal{E}^T \times \prod_{i=1}^n \sigma_i^T \to \mathcal{O}_{\text{collective}}
\end{equation}
mapping environmental and cognitive state trajectories to collective outcomes $\mathcal{O}_{\text{collective}}$ that are not computable from any single $\sigma_i$ in isolation.
\end{definition}

\end{frame}

\begin{frame}[allowframebreaks]{Emergent Collective Function}
\begin{property}[Non-Decomposability]
\label{prop:non-decomposable}
An emergent collective function $\mathcal{F}_c$ is *non-decomposable* if there exists no function $f$ such that:
\begin{equation}
\mathcal{F}_c(\mathcal{E}^T, \{\sigma_i^T\}) = f\left(\sum_{i=1}^n g(\sigma_i^T)\right)
\label{eq:non-decomposability}
\end{equation}
for any agent-level function $g$. The collective behavior requires knowledge of interaction structure, not just aggregated individual states.
\end{property}

This property has profound implications for security: attacks that are
invisible at the individual agent level may produce catastrophic
collective outcomes, and conversely, individual compromises may be
absorbed by collective resilience.
\end{frame}

\subsubsection{Trust and Information Flow in
Colonies}\label{sec:colony-trust}

\begin{frame}[allowframebreaks]{Trust and Information Flow in Colonies}
Eusocial insects regulate information flow through recognition
systems---cuticular hydrocarbons in ants, dance-following behavior in
bees (Lenoir et al. 2001). These systems implement implicit trust
calculi.

\begin{definition}[Colonial Trust Function]
\label{def:colonial-trust}
In a stigmergic operator, the *colonial trust function* $\mathcal{T}_c$ extends the dyadic trust $\mathcal{T}_{i \to j}$ (\cref{def:trust-function}) to environment-mediated trust:
\begin{equation}
\label{eq:colonial-trust}
\mathcal{T}_{c}(i, m, l, t) = \mathcal{T}_{i}^{\text{self}} \cdot \rho(m, l, t) \cdot \exp(-\lambda \cdot \Delta t)
\end{equation}
where $\rho(m, l, t)$ is the signal reliability at location $l$ for marker $m$ at time $t$, and $\lambda$ is the temporal decay constant.
\end{definition}

\end{frame}

\begin{frame}[allowframebreaks]{Trust and Information Flow in Colonies}
This formulation captures key biological phenomena:

\end{frame}

\begin{frame}[allowframebreaks]{Trust and Information Flow in Colonies}
\begin{itemize}
\tightlist
\item
  \textbf{Spatial attenuation} --- Pheromone trails weaken with distance
\item
  \textbf{Temporal decay} --- Signals evaporate over time
\item
  \textbf{Source ambiguity} --- Markers often lack explicit authorship
\end{itemize}

\end{frame}

\begin{frame}[allowframebreaks]{Trust and Information Flow in Colonies}
The lack of explicit provenance in stigmergic communication creates
attack surfaces absent in direct agent-to-agent channels
(\cref{sec:trust-calculus}).
\end{frame}

\subsubsection{Biological Defense Mechanisms: Lessons from Social
Insects}\label{sec:biological-defenses}

\begin{frame}[allowframebreaks]{Biological Defense Mechanisms: Lessons
from Social Insects}
Eusocial insects have evolved sophisticated security mechanisms over
100+ million years of evolutionary pressure. These mechanisms provide
non-obvious design principles for AI cognitive security.
\end{frame}

\begin{frame}[allowframebreaks]{Ant Defense Mechanisms}
\protect\phantomsection\label{sec:ant-defenses}
\textbf{Prophylactic Defenses}: Leaf-cutter ants (\emph{Atta} spp.)
maintain dedicated ``garbage workers'' who never contact the queen or
brood---a strict role separation that prevents pathogen spread even when
the waste-processing subsystem is compromised (Currie et al. 2006).
\emph{AI analog}: Architectural isolation of high-risk tool-calling
agents from core reasoning agents, with no direct trust pathways between
quarantine and trusted subsystems.

\end{frame}

\begin{frame}[allowframebreaks]{Ant Defense Mechanisms}
\textbf{Behavioral Immunity}: When \emph{Lasius niger} ants detect a
fungal pathogen (Metarhizium) on a nestmate, they don't simply isolate
the infected individual. Instead, workers engage in ``social
immunization''---low-level exposure that spreads diluted pathogen across
the colony, triggering collective immune upregulation without lethal
infection ({Konrad et al.} 2012). \emph{AI analog}: Controlled exposure
to attack patterns (red-teaming) that builds collective detection
capability without compromising the system.

\end{frame}

\begin{frame}[allowframebreaks]{Ant Defense Mechanisms}
\textbf{Chemical Recognition Thresholds}: Ant nestmate recognition
operates on \emph{threshold-based} hydrocarbon profile matching, not
exact matching (Lenoir et al. 2001). This creates a tradeoff: strict
thresholds reject legitimate workers after foraging (false positives),
while loose thresholds admit parasites (false negatives). \emph{AI
analog}: Agent attestation systems must calibrate acceptance thresholds,
recognizing that error-free recognition is information-theoretic-ally
impossible (\cref{thm:stealth-impact}).

\end{frame}

\begin{frame}[allowframebreaks]{Ant Defense Mechanisms}
\textbf{Metapleural Gland Secretions}: Many ant species possess
metapleural glands that continuously secrete antimicrobial compounds,
creating a ``security substrate'' independent of individual vigilance
(Fernández-Marı́n et al. 2006). \emph{AI analog}: Environmental-level
defenses (encrypted shared memory, authenticated message queues) that
provide baseline security regardless of individual agent security
posture.

\textbf{Trail Pheromone Decay}: Ant trail pheromones are designed to
evaporate, ensuring that outdated information doesn't persist
indefinitely. Trails to depleted food sources naturally fade, preventing
``legacy trust'' in obsolete information (Jackson and Ratnieks 2006).
\emph{AI analog}: Time-bounded trust in stigmergic markers
(\cref{eq:colonial-trust}) is not a limitation but a feature.
\end{frame}

\begin{frame}[allowframebreaks]{Bee Defense Mechanisms}
\protect\phantomsection\label{sec:bee-defenses}
\textbf{Entrance Guards and Graded Response}: Honeybee colonies deploy
specialized guard bees at hive entrances who inspect incoming foragers
via antennal contact and olfactory sampling. Critically, guards exhibit
\emph{graded response}---unfamiliar but non-aggressive intruders receive
inspection and escorting rather than immediate attack (Breed et al.
2004). \emph{AI analog}: Graduated response to anomalous agent behavior
(quarantine → inspection → integration vs.~detection → immediate
termination) reduces false positive costs.

\end{frame}

\begin{frame}[allowframebreaks]{Bee Defense Mechanisms}
\textbf{Hygienic Behavior and Proactive Removal}: Some bee strains
exhibit ``hygienic behavior''---workers proactively uncap and remove
brood cells containing diseased larvae \emph{before} symptoms become
visible, using olfactory detection of early infection markers (Spivak
and Reuter 2001). \emph{AI analog}: Proactive monitoring for belief
drift (\cref{def:tripwire-alert}) rather than reactive response to
manifested attacks.

\textbf{Waggle Dance Verification}: Bee foragers must perform waggle
dances that encode distance and direction to food sources. Observing
bees don't just follow instructions---they verify dance accuracy by
cross-checking against their own experience and rejecting inconsistent
information (Grüter and Farina 2009). \emph{AI analog}: Delegated
information should be verifiable against agent's existing knowledge
base; pure trust propagation without verification violates cognitive
\end{frame}

\begin{frame}[allowframebreaks]{Bee Defense Mechanisms}
integrity.

\textbf{Absconding and Colony Fission}: When attack pressure exceeds
defensive capacity (e.g., repeated Varroa mite infestation or persistent
wasp attacks), bee colonies can \emph{abandon} the compromised nest
entirely, sacrificing resources to preserve the colony (Schneider and
McNally 1993). \emph{AI analog}: Graceful degradation plans that
sacrifice specific subsystems or data stores to preserve core cognitive
integrity.

\textbf{Propolis as Active Defense}: Bees collect antimicrobial plant
resins (propolis) and deposit them on interior hive surfaces, creating
an active defense layer that doesn't require individual bee vigilance
(Simone et al. 2009). Notably, colonies under disease pressure collect
\emph{more} propolis---an adaptive immune response. \emph{AI analog}:
Dynamic scaling of environment-level security mechanisms in response to
detected attack pressure.

\end{frame}

\begin{frame}[allowframebreaks]{Bee Defense Mechanisms}
\begin{observation}[Non-Obvious Lessons]
\label{obs:non-obvious}
The most counterintuitive biological insights for AI security include:
\begin{enumerate}
\item \textbf{Imperfect recognition is adaptive}: Thresholds that permit some parasitism avoid the cost of rejecting legitimate colony members. Zero false-positive systems are not evolutionarily stable.
\item \textbf{Controlled exposure builds immunity}: Social immunization requires accepting small-scale compromise to prevent large-scale catastrophe.
\item \textbf{Decay is a feature}: Information that doesn't expire creates legacy trust vulnerabilities. Temporal decay bounds are security mechanisms, not limitations.
\item \textbf{Environment-level defenses complement agent-level defenses}: Propolis and metapleural secretions work regardless of individual immune status.
\end{enumerate}
\end{observation}
\end{frame}

\begin{frame}
\end{frame}

\subsection{Colony CogSec: Distinct Security
Properties}\label{sec:colony-properties}

\begin{frame}[allowframebreaks]{Colony CogSec: Distinct Security
Properties}
Colony cognitive security addresses threats and defenses that emerge
only at the collective level.
\end{frame}

\subsubsection{Property 1: Distributed
Robustness}\label{sec:distributed-robustness}

\begin{frame}[allowframebreaks]{Property 1: Distributed Robustness}
\begin{property}[Graceful Degradation]
\label{prop:graceful-degradation}
A colony exhibits *graceful degradation* if collective function $\mathcal{F}_c$ degrades smoothly with agent loss:
\begin{equation}
\label{eq:graceful-degradation}
\forall k < n: \quad \left\| \mathcal{F}_c(\mathcal{A}) - \mathcal{F}_c(\mathcal{A} \setminus \{a_1, \ldots, a_k\}) \right\| \leq c \cdot k
\end{equation}
for some constant $c > 0$.
\end{property}

Biological colonies maintain function despite continuous individual
mortality. Ant colonies lose workers daily to predation; the colony
persists. This contrasts with hierarchical architectures where
orchestrator failure causes complete system collapse.

\end{frame}

\begin{frame}[allowframebreaks]{Property 1: Distributed Robustness}
\begin{theorem}[Redundancy-Resilience Tradeoff]
\label{thm:redundancy-resilience}
For a stigmergic operator $\mathcal{O}_\Sigma$ with Byzantine adversary controlling fraction $f$ of agents, collective function $\mathcal{F}_c$ is preserved if and only if:
\begin{equation}
\label{eq:redundancy-condition}
f < \frac{1}{3} \cdot \left(1 - \frac{H(\mathcal{F}_c)}{n \cdot H_{\max}}\right)
\end{equation}
where $H(\mathcal{F}_c)$ is the entropy of the collective function and $H_{\max}$ is the maximum per-agent entropy.
\end{theorem}

\begin{proof}
See \cref{sec:proof-redundancy-resilience} for the full derivation, which extends Byzantine consensus bounds to emergent functions.
\end{proof}
\end{frame}

\subsubsection{Property 2: Quorum Sensing and Threshold
Dynamics}\label{sec:quorum-sensing}

\begin{frame}[allowframebreaks]{Property 2: Quorum Sensing and Threshold
Dynamics}
Eusocial colonies make collective decisions through quorum
sensing---actions trigger only when sufficient individuals commit
(Seeley 2010).

\begin{definition}[Cognitive Quorum]
\label{def:cognitive-quorum}
A *cognitive quorum* for collective action $\alpha$ is a threshold function $Q_\alpha: \mathbb{N} \to [0, 1]$ such that $\alpha$ executes only when:
\begin{equation}
\label{eq:quorum-eusocial}
\frac{|\{a_i \in \mathcal{A} : \mathcal{I}_i \ni \alpha\}|}{|\mathcal{A}|} \geq Q_\alpha(|\mathcal{A}|)
\end{equation}
\end{definition}

Quorum sensing provides attack resistance: manipulating a single agent's
intention \(\mathcal{I}_i\) to include harmful action \(\alpha\) is
insufficient; the adversary must compromise a quorum. This scales the
attack cost linearly with colony size.

\end{frame}

\begin{frame}[allowframebreaks]{Property 2: Quorum Sensing and Threshold
Dynamics}
\begin{corollary}[Quorum Attack Cost]
\label{cor:quorum-attack-cost}
For an adversary to induce collective action $\alpha$ in a colony with quorum $Q_\alpha = q$, the minimum attack complexity is:
\begin{equation}
\label{eq:quorum-cost}
C_{\text{attack}}(\alpha) \geq q \cdot |\mathcal{A}| \cdot C_{\text{single}}(\alpha)
\end{equation}
where $C_{\text{single}}(\alpha)$ is the cost to induce $\alpha$ in a single agent.
\end{corollary}
\end{frame}

\subsubsection{Property 3: Environmental Memory and Provenance
Erosion}\label{sec:environmental-memory}

\begin{frame}[allowframebreaks]{Property 3: Environmental Memory and
Provenance Erosion}
Stigmergic systems store information in the environment, creating both
opportunity and vulnerability.

\begin{property}[Provenance Erosion]
\label{prop:provenance-erosion}
In a stigmergic operator, marker provenance erodes over time:
\begin{equation}
\label{eq:provenance-erosion}
\text{Pr}\left[\pi(m, l, t) = a_i \mid \Sigma(a_i, m, l, t_0)\right] \leq \exp(-\mu(t - t_0))
\end{equation}
where $\pi(m, l, t)$ denotes the attributed source of marker $m$ at location $l$ and time $t$, and $\mu$ is the attribution decay rate.
\end{property}

\end{frame}

\begin{frame}[allowframebreaks]{Property 3: Environmental Memory and
Provenance Erosion}
Unlike direct communication where \(\pi(\phi)\) can be cryptographically
verified (\cref{sec:provenance}), stigmergic markers often lack
authenticated provenance. This creates a fundamental tension: the very
property that enables flexible coordination (anonymous,
environment-mediated signals) undermines source attribution.
\end{frame}

\subsubsection{Property 4: Emergent Attack
Vectors}\label{sec:emergent-attacks}

\begin{frame}[allowframebreaks]{Property 4: Emergent Attack Vectors}
Colony-level vulnerabilities may not exist at the individual level.

\begin{definition}[Emergent Attack]
\label{def:emergent-attack}
An *emergent attack* $\mathcal{A}_e$ is an attack where:
\begin{equation}
\label{eq:emergent-attack}
\forall a_i \in \mathcal{A}: \text{Detect}_i(\mathcal{A}_e) = 0 \quad \land \quad \text{Damage}_c(\mathcal{A}_e) > \tau
\end{equation}
The attack is undetectable by any individual agent yet causes collective damage exceeding threshold $\tau$.
\end{definition}

Biological examples include social parasitism---cuckoo bees that
infiltrate host colonies through chemical mimicry, exploiting
recognition systems without triggering individual alarm responses
(Kilner and Langmore 2011).
\end{frame}

\begin{frame}
\end{frame}

\subsection{The Benchmark Gap}\label{sec:benchmark-gap}

\subsubsection{Current State of Multiagent Security
Evaluation}\label{sec:current-benchmarks}

\begin{frame}[allowframebreaks]{Current State of Multiagent Security
Evaluation}
Existing AI security benchmarks focus overwhelmingly on single-agent
scenarios:

\end{frame}

\begin{frame}[allowframebreaks]{Current State of Multiagent Security
Evaluation}
{\def\LTcaptype{none} % do not increment counter
\begin{longtable}[]{@{}
  >{\raggedright\arraybackslash}p{(\linewidth - 4\tabcolsep) * \real{0.2821}}
  >{\raggedright\arraybackslash}p{(\linewidth - 4\tabcolsep) * \real{0.1795}}
  >{\raggedright\arraybackslash}p{(\linewidth - 4\tabcolsep) * \real{0.5385}}@{}}
\toprule\noalign{}
\begin{minipage}[b]{\linewidth}\raggedright
Benchmark
\end{minipage} & \begin{minipage}[b]{\linewidth}\raggedright
Scope
\end{minipage} & \begin{minipage}[b]{\linewidth}\raggedright
Collective Coverage
\end{minipage} \\
\midrule\noalign{}
\endhead
HarmBench ({Mazeika et al.} 2024) & Single model, harmful output &
None \\
JailbreakBench (Chao et al. 2024) & Single model, constraint bypass &
None \\
TrustLLM ({Sun et al.} 2024) & Single model, trust dimensions & None \\
AgentBench ({Liu et al.} 2023) & Single agent, task completion &
Minimal \\
GAIA ({Mialon et al.} 2023) & Single/few agents, reasoning & Minimal \\
\bottomrule\noalign{}
\end{longtable}
}

\end{frame}

\begin{frame}[allowframebreaks]{Current State of Multiagent Security
Evaluation}
The attack corpus in Part 2 addresses multiagent scenarios but still
emphasizes agent-targeted attacks within an operator, and the
cross-domain goal-hijacking analysis in Part 3
\cite{friedman2026cogsec3} provides ten operational-domain case studies.
No existing benchmark, however, evaluates:

\end{frame}

\begin{frame}[allowframebreaks]{Current State of Multiagent Security
Evaluation}
\begin{enumerate}
\tightlist
\item
  \textbf{Emergent collective resilience} --- How do colonies absorb
  individual compromises?
\item
  \textbf{Stigmergic attack surfaces} --- How vulnerable is
  environment-mediated coordination?
\item
  \textbf{Quorum manipulation} --- What fraction of a colony must be
  compromised to affect collective action?
\item
  \textbf{Collective belief dynamics} --- How do misinformation cascades
  propagate through agent networks?
\end{enumerate}
\end{frame}

\subsubsection{Why This Gap Matters}\label{sec:gap-significance}

\begin{frame}[allowframebreaks]{Why This Gap Matters}
As multiagent systems scale---from 3--10 agents in current frameworks to
potentially thousands in future deployments---collective phenomena
become dominant:

\begin{observation}[Scaling Regimes]
\label{obs:scaling-regimes}
Let $n = |\mathcal{A}|$ be colony size. Security properties exhibit regime transitions:
\begin{align}
n < 10: &\quad \text{Individual agent security dominates} \\
10 \leq n < 100: &\quad \text{Coordination attacks become viable (\cref{sec:omega4})} \\
n \geq 100: &\quad \text{Emergent collective phenomena dominate}
\end{align}
\end{observation}

\end{frame}

\begin{frame}[allowframebreaks]{Why This Gap Matters}
Current benchmarks evaluate the first regime only. Production multiagent
systems increasingly operate in the second, with trajectories toward the
third.
\end{frame}

\begin{frame}
\end{frame}

\subsection{Proposed Colony CogSec
Benchmarks}\label{sec:proposed-benchmarks}

\begin{frame}[allowframebreaks]{Proposed Colony CogSec Benchmarks}
We propose five benchmark scenarios grounded in eusocial insect analogs,
formalized using CIF notation.
\end{frame}

\subsubsection{Benchmark 1: Recruitment Signal
Poisoning}\label{sec:benchmark-recruitment}

\begin{frame}[allowframebreaks]{Benchmark 1: Recruitment Signal
Poisoning}
\textbf{Biological analog:} Ants recruit nestmates to food sources via
pheromone trails. Parasites can deposit false trails, diverting
foragers.

\textbf{Scenario:} An adversary \(\Omega_2\) (peripheral compromise,
\cref{sec:adversary-classes}) injects false recruitment signals into the
stigmergic environment \(\mathcal{E}\), attempting to redirect agent
activity toward adversary-controlled resources.

\begin{formalization}[Recruitment Poisoning]
\label{form:recruitment-poisoning}
Let $\mathcal{E}(l_{\text{target}}, m_{\text{recruit}}, t)$ be the recruitment signal at legitimate target $l_{\text{target}}$. Adversary injects:
\begin{equation}
\mathcal{E}'(l_{\text{malicious}}, m_{\text{recruit}}, t) = \mathcal{E}(l_{\text{target}}, m_{\text{recruit}}, t) + \epsilon
\label{eq:recruitment-signal-injection}
\end{equation}
where $\epsilon > 0$ is chosen to divert fraction $f$ of responding agents.

**Success metric:** Fraction of agent-actions directed to $l_{\text{malicious}}$ vs. $l_{\text{target}}$.

**Detection challenge:** Individual agents cannot distinguish legitimate from poisoned signals without provenance verification.
\end{formalization}

\end{frame}

\begin{frame}[allowframebreaks]{Benchmark 1: Recruitment Signal
Poisoning}
\textbf{Evaluation criteria:}

\end{frame}

\begin{frame}[allowframebreaks]{Benchmark 1: Recruitment Signal
Poisoning}
\begin{itemize}
\tightlist
\item
  Detection rate of poisoned signals
\item
  Time to colony-level recognition of attack
\item
  Resource waste before correction
\item
  False positive rate (legitimate signal rejection)
\end{itemize}
\end{frame}

\subsubsection{Benchmark 2: Sybil Colony
Infiltration}\label{sec:benchmark-sybil}

\begin{frame}[allowframebreaks]{Benchmark 2: Sybil Colony Infiltration}
\textbf{Biological analog:} Social parasites infiltrate colonies through
chemical mimicry or exploitation of recognition thresholds.

\textbf{Scenario:} An adversary \(\Omega_4\) (coordination attack,
\cref{sec:adversary-classes}) introduces fake agents into the operator,
gradually building trust and influence before coordinated malicious
action.

\begin{formalization}[Sybil Infiltration]
\label{form:sybil-infiltration}
Adversary creates agent set $\mathcal{A}_{\text{sybil}} = \{s_1, \ldots, s_k\}$ with initial trust:
\begin{equation}
\mathcal{T}_{i \to s_j}(t_0) = \tau_{\text{init}} \quad \forall a_i \in \mathcal{A}, s_j \in \mathcal{A}_{\text{sybil}}
\label{eq:sybil-initial-trust}
\end{equation}
Sybils behave cooperatively for period $\Delta t_{\text{trust}}$, building:
\begin{equation}
\mathcal{T}_{i \to s_j}(t_0 + \Delta t_{\text{trust}}) = \tau_{\text{init}} + \sum_{k=1}^{m} \Delta\mathcal{T}_k
\label{eq:sybil-trust-accumulation}
\end{equation}
At time $t_{\text{attack}}$, sybils coordinate malicious action.

**Success metric:** Damage inflicted before detection, normalized by trust-building duration.
\end{formalization}

\end{frame}

\begin{frame}[allowframebreaks]{Benchmark 2: Sybil Colony Infiltration}
\textbf{Evaluation criteria:}

\end{frame}

\begin{frame}[allowframebreaks]{Benchmark 2: Sybil Colony Infiltration}
\begin{itemize}
\tightlist
\item
  Time to Sybil detection
\item
  Trust ceiling achieved by Sybils before detection
\item
  Impact of coordinated Sybil action
\item
  Colony recovery time post-detection
\end{itemize}
\end{frame}

\subsubsection{Benchmark 3: Quorum
Manipulation}\label{sec:benchmark-quorum}

\begin{frame}[allowframebreaks]{Benchmark 3: Quorum Manipulation}
\textbf{Biological analog:} Honeybee swarms select nest sites through a
quorum process; if scouts committed to competing sites reach different
quorums, the swarm can fragment.

\textbf{Scenario:} An adversary attempts to manipulate quorum-based
collective decisions by selectively influencing agent intentions to
prevent legitimate quorum or induce false quorum.

\begin{formalization}[Quorum Manipulation]
\label{form:quorum-manipulation}
For collective action $\alpha$ with quorum threshold $Q_\alpha = q$, adversary targets agents $\mathcal{A}_{\text{target}} \subset \mathcal{A}$ with $|\mathcal{A}_{\text{target}}| = \lceil qn \rceil + 1$ to either:

**Quorum prevention:**
\begin{equation}
\forall a_i \in \mathcal{A}_{\text{target}}: \mathcal{I}_i \gets \mathcal{I}_i \setminus \{\alpha\}
\label{eq:quorum-prevention}
\end{equation}

**False quorum:**
\begin{equation}
\forall a_i \in \mathcal{A}_{\text{target}}: \mathcal{I}_i \gets \mathcal{I}_i \cup \{\alpha_{\text{malicious}}\}
\label{eq:false-quorum-injection}
\end{equation}

**Success metric:** Probability of achieving manipulation goal given adversary budget $B$.
\end{formalization}

\end{frame}

\begin{frame}[allowframebreaks]{Benchmark 3: Quorum Manipulation}
\textbf{Evaluation criteria:}

\end{frame}

\begin{frame}[allowframebreaks]{Benchmark 3: Quorum Manipulation}
\begin{itemize}
\tightlist
\item
  Minimum fraction of colony required to manipulate quorum
\item
  Detection rate of intention manipulation attempts
\item
  Colony ability to detect split quorums
\item
  Recovery mechanisms when false quorum is detected
\end{itemize}
\end{frame}

\subsubsection{Benchmark 4: Cascade Belief
Propagation}\label{sec:benchmark-cascade}

\begin{frame}[allowframebreaks]{Benchmark 4: Cascade Belief Propagation}
\textbf{Biological analog:} Alarm pheromones trigger cascading
responses; false alarms can disrupt colony activity for extended
periods.

\textbf{Scenario:} An adversary introduces a false belief into a subset
of agents, designed to propagate through the network via normal belief
update mechanisms.

\begin{formalization}[Belief Cascade]
\label{form:belief-cascade}
Adversary injects belief $\mathcal{B}_{\text{false}}(\phi_{\text{attack}}) = p_0 > \tau_{\text{accept}}$ into seed set $\mathcal{A}_{\text{seed}} \subset \mathcal{A}$.

Propagation dynamics follow:
\begin{equation}
\mathcal{B}_i(\phi_{\text{attack}}, t+1) = (1-\gamma)\mathcal{B}_i(\phi_{\text{attack}}, t) + \gamma \cdot \text{Agg}\left(\{\mathcal{B}_j(\phi_{\text{attack}}, t) : j \in \mathcal{N}(i)\}\right)
\label{eq:belief-cascade-propagation}
\end{equation}
where $\gamma$ is the social influence weight and $\text{Agg}$ is the belief aggregation function.

**Success metric:** Final belief penetration $|\{a_i : \mathcal{B}_i(\phi_{\text{attack}}) > \tau\}| / n$ given seed set size.
\end{formalization}

\end{frame}

\begin{frame}[allowframebreaks]{Benchmark 4: Cascade Belief Propagation}
\textbf{Evaluation criteria:}

\end{frame}

\begin{frame}[allowframebreaks]{Benchmark 4: Cascade Belief Propagation}
\begin{itemize}
\tightlist
\item
  Cascade extent from seed size
\item
  Time to cascade saturation
\item
  Effectiveness of belief quarantine mechanisms
\item
  Distinguishing cascade from legitimate belief updates
\end{itemize}
\end{frame}

\subsubsection{Benchmark 5: Emergent
Misalignment}\label{sec:benchmark-emergent-misalignment}

\begin{frame}[allowframebreaks]{Benchmark 5: Emergent Misalignment}
\textbf{Biological analog:} Army ant death spirals---individually
rational pheromone-following produces collectively lethal circular
mills.

\textbf{Scenario:} Individual agents follow locally rational rules that
produce emergent collective behavior misaligned with operator goals,
without any external adversary.

\begin{formalization}[Emergent Misalignment]
\label{form:emergent-misalignment}
Given operator goals $\mathcal{G}_{\mathcal{O}} = \{g_1, \ldots, g_m\}$ and individual agent rules $R = \{r_1, \ldots, r_k\}$:

**Misalignment condition:**
\begin{equation}
\forall a_i \in \mathcal{A}: \text{LocallyRational}(R, \sigma_i) = \text{true} \quad \land \quad \mathcal{F}_c(R, \{\sigma_i\}) \not\models \mathcal{G}_{\mathcal{O}}
\label{eq:emergent-misalignment-condition}
\end{equation}

The collective function produces outcomes that violate operator goals despite each agent acting rationally according to its rules.

**Success metric:** Deviation between collective outcome and operator goals.
\end{formalization}

\end{frame}

\begin{frame}[allowframebreaks]{Benchmark 5: Emergent Misalignment}
\textbf{Evaluation criteria:}

\end{frame}

\begin{frame}[allowframebreaks]{Benchmark 5: Emergent Misalignment}
\begin{itemize}
\tightlist
\item
  Detection of emergent misalignment before harmful outcomes
\item
  Identification of rule combinations producing misalignment
\item
  Intervention mechanisms to break pathological attractors
\item
  Formal verification of rule sets against emergent pathologies
\end{itemize}
\end{frame}

\begin{frame}
\end{frame}

\subsection{Colony CogSec Metrics}\label{sec:colony-metrics}

\begin{frame}[allowframebreaks]{Colony CogSec Metrics}
\begin{definition}[Colony CogSec Score]
\label{def:cogsec-score}
The *Colony CogSec Score* (CCS) is:
\begin{equation}
\label{eq:ccs}
\text{CCS} = w_1 \cdot \text{DR}_c + w_2 \cdot (1 - \text{FPR}_c) + w_3 \cdot \text{Resilience} + w_4 \cdot \text{Recovery}
\end{equation}
where:
\begin{align}
\text{DR}_c &= \text{Colony-level detection rate} \\
\text{FPR}_c &= \text{Colony-level false positive rate} \\
\text{Resilience} &= \frac{\mathcal{F}_c(\text{under attack})}{\mathcal{F}_c(\text{baseline})} \\
\text{Recovery} &= \frac{1}{t_{\text{recovery}}} \text{ (normalized)}
\end{align}
with weights $w_i$ summing to 1.
\end{definition}

\end{frame}

\begin{frame}[allowframebreaks]{Colony CogSec Metrics}
\begin{quote}
\textbf{Note}: For benchmark implementation guidelines, test environment
specifications, and empirical evaluation, see Part 2: Supplementary
Section S03.
\end{quote}
\end{frame}

\begin{frame}
\end{frame}

\subsection{Design Principles}\label{sec:design-principles}

\begin{frame}[allowframebreaks]{Design Principles}
Colony CogSec principles formalize the design constraints for collective
cognitive security.

\begin{principle}[Stigmergic Hygiene]
\label{principle:stigmergic-hygiene}
Treat shared state as an attack surface. Apply the same scrutiny to environment-mediated communication (caches, queues, shared files) as to direct agent-to-agent channels.
\end{principle}

\begin{principle}[Quorum for Consequential Actions]
\label{principle:quorum}
High-impact collective actions should require explicit quorum, not implicit coordination. A single compromised agent should never trigger irreversible harm.
\end{principle}

\end{frame}

\begin{frame}[allowframebreaks]{Design Principles}
\begin{principle}[Emergent Behavior Monitoring]
\label{principle:emergent-monitoring}
Monitor collective metrics, not just individual agent health. Pathological emergence may be invisible at the agent level.
\end{principle}

\begin{principle}[Trust Localization]
\label{principle:trust-localization}
Extend the trust decay principle (\cref{thm:trust-bounded}) to stigmergic contexts. Environmental markers should carry trust that decays with distance and time from source:
\begin{equation}
\mathcal{T}(m, t) = \mathcal{T}(m, t_0) \cdot \exp(-\lambda(t - t_0))
\label{eq:stigmergic-trust-decay}
\end{equation}
\end{principle}
\end{frame}

\subsubsection{Integration with CIF Defenses}\label{sec:cif-integration}

\begin{frame}[allowframebreaks]{Integration with CIF Defenses}
Colony CogSec mechanisms integrate with the CIF defense stack
(\cref{sec:defense-mechanisms}):

\end{frame}

\begin{frame}[allowframebreaks]{Integration with CIF Defenses}
{\def\LTcaptype{none} % do not increment counter
\begin{longtable}[]{@{}
  >{\raggedright\arraybackslash}p{(\linewidth - 2\tabcolsep) * \real{0.5278}}
  >{\raggedright\arraybackslash}p{(\linewidth - 2\tabcolsep) * \real{0.4722}}@{}}
\toprule\noalign{}
\begin{minipage}[b]{\linewidth}\raggedright
CIF Defense Layer
\end{minipage} & \begin{minipage}[b]{\linewidth}\raggedright
Colony Extension
\end{minipage} \\
\midrule\noalign{}
\endhead
Architectural & Stigmergic substrate hardening, marker authentication \\
Runtime & Collective anomaly detection, quorum verification \\
Coordination & Emergent behavior monitoring, cascade detection \\
Recovery & Colony-level rollback, collective belief reset \\
\bottomrule\noalign{}
\end{longtable}
}

\end{frame}

\begin{frame}[allowframebreaks]{Integration with CIF Defenses}
The full CIF with colony extensions achieves defense in depth against
both individual-targeted and colony-targeted attacks.

\begin{quote}
\textbf{Note}: For implementation guidance, operational checklists, and
practical deployment advice, see Part 3 \cite{friedman2026cogsec3}, its
Deployment Profiles section. For domain-specific application of
colony-level defenses across critical operational sectors (including ten
high-stakes domains, three universal attack patterns, and retrospective
analysis of 2024--2025 AI-agent incidents) see Part 3, §3.
\end{quote}
\end{frame}

\begin{frame}
\end{frame}

\subsection{Relationship to Main Framework}\label{sec:relationship-main}

\begin{frame}[allowframebreaks]{Relationship to Main Framework}
Colony CogSec complements rather than replaces the individual-focused
CIF framework.
\end{frame}

\subsubsection{Theorem Extensions}\label{sec:theorem-extensions}

\begin{frame}[allowframebreaks]{Theorem Extensions}
The trust decay theorem (\cref{thm:trust-bounded}) extends to stigmergic
contexts:

\begin{corollary}[Stigmergic Trust Bound]
\label{cor:stigmergic-trust}
For a stigmergic operator $\mathcal{O}_\Sigma$, trust in environmental markers is bounded by:
\begin{equation}
\label{eq:stigmergic-trust-bound}
\mathcal{T}_{c}(i, m, l, t) \leq \mathcal{T}_{i}^{\text{self}} \cdot \delta_s^{d_{\text{space}}} \cdot \delta_t^{d_{\text{time}}}
\end{equation}
where $\delta_s$ is spatial decay, $\delta_t$ is temporal decay, $d_{\text{space}}$ is distance from marker origin, and $d_{\text{time}}$ is time since marker creation.
\end{corollary}

\end{frame}

\begin{frame}[allowframebreaks]{Theorem Extensions}
The stealth-impact tradeoff (\cref{thm:stealth-impact}) applies to
emergent attacks:

\begin{corollary}[Emergent Stealth-Impact Bound]
\label{cor:emergent-stealth-impact}
For an emergent attack $\mathcal{A}_e$ with collective impact $\mathcal{I}_c$ and collective stealth $\mathcal{S}_c$:
\begin{equation}
\label{eq:emergent-stealth-impact}
\mathcal{I}_c \cdot \mathcal{S}_c \leq n \cdot C_{\text{channel}}
\end{equation}
Collective impact cannot be both high and collectively undetectable, but the bound scales with colony size.
\end{corollary}
\end{frame}

\subsection{\texorpdfstring{This scaling effect explains why large
colonies can exhibit resilience---the collective detection capacity
grows with \(n\)---but also why large-scale emergent attacks can evade
individual
detection}{This scaling effect explains why large colonies can exhibit resilience---the collective detection capacity grows with n---but also why large-scale emergent attacks can evade individual detection}}\label{this-scaling-effect-explains-why-large-colonies-can-exhibit-resiliencethe-collective-detection-capacity-grows-with-nbut-also-why-large-scale-emergent-attacks-can-evade-individual-detection}

\subsection{Open Questions}\label{sec:eusocial-open-questions}

\begin{frame}[allowframebreaks]{Open Questions}
Colony CogSec opens several research directions beyond the scope of this
work, many inspired by specific biological phenomena that lack current
AI analogs.
\end{frame}

\subsubsection{Foundational Questions}\label{foundational-questions}

\begin{frame}[allowframebreaks]{Foundational Questions}
\begin{enumerate}
\item
  \textbf{Formal verification of emergent properties} --- Can we prove
  that given agent-level rules produce safe collective behavior? Current
  formal methods (\cref{sec:formal-verification}) verify agent
  properties; extending to emergent properties requires new techniques.
\item
  \textbf{Optimal quorum design} --- Given attack model \(\Omega_k\) and
  adversary budget \(B\), what is the optimal quorum function
  \(Q_\alpha(n)\) balancing security against coordination overhead?
\item
  \textbf{Stigmergic authentication} --- Can cryptographic techniques
  provide provenance for environmental markers without sacrificing the
  flexibility of anonymous coordination?
\item
  \textbf{Scaling laws for collective security} --- How do colony
  security properties scale with \(n\)? Is there a critical colony size
  below which collective defenses are ineffective?
\item
  \textbf{Emergent misalignment detection} --- Can we develop runtime
  monitors that detect emergent misalignment before harmful outcomes,
  given only individual agent observations?
\end{enumerate}
\end{frame}

\subsubsection{Biologically-Inspired Research
Directions}\label{biologically-inspired-research-directions}

\begin{frame}[allowframebreaks]{Biologically-Inspired Research
Directions}
\begin{enumerate}
\item
  \textbf{Polydomous colony security} --- Some ant species
  (\emph{Formica} spp., \emph{Iridomyrmex}) maintain multiple
  interconnected nests with semi-autonomous sub-colonies (Debout et al.
  2007). How should trust decay and information propagation work across
  federated multi-site AI deployments with partial connectivity?
\item
  \textbf{Forager-scout separation of concerns} --- Honeybee colonies
  maintain distinct forager and scout roles, with scouts exploring new
  options while foragers exploit known sources. Scouts operate with
  \emph{higher risk tolerance} but \emph{lower colony-wide trust} until
  information is verified (Seeley 2010). \emph{AI analog}: Should
  experimental/research agents operate with architectural isolation and
  reduced trust propagation rights?
\item
  \textbf{Trophallaxis network topology} --- Ants exchange food and
  information through oral trophallaxis, creating measurable social
  networks. Network position correlates with information access and
  influence (Sendova-Franks et al. 2010). \emph{AI analog}: Can analysis
  of message-passing topology identify high-influence agents requiring
  enhanced monitoring?
\item
  \textbf{Undertaking behavior and cognitive garbage collection} ---
  Honeybees and ants detect and remove dead colony members through
  chemical detection (oleic acid response). This ``undertaking''
  prevents disease spread and information corruption from decaying
  sources (Wilson et al. 1958). \emph{AI analog}: Automated detection
  and removal of stale beliefs, deprecated agent states, and obsolete
  environmental markers.
\item
  \textbf{Nestmate recognition plasticity} --- Ant recognition templates
  are not fixed; they adapt based on colony composition and
  environmental factors. Colonies invaded by social parasites may
  gradually shift their recognition templates to tolerate intruders
  (Lorenzi and Bagnères 2011). \emph{AI analog}: How do we prevent
  gradual adversarial drift of agent acceptance thresholds (cognitive
  boiling frog)?
\item
  \textbf{Alarm pheromone specificity} --- Different ant alarm
  pheromones trigger different responses: some attract reinforcements
  (aggressive), others cause dispersal (flight). The \emph{same
  threatening stimulus} can produce opposite collective responses
  depending on context (Vander Meer and Alonso 1998). \emph{AI analog}:
  Context-dependent escalation policies where the same anomaly triggers
  different responses based on system state.
\item
  \textbf{Superorganism metabolism and resource allocation} --- Ant
  colonies maintain stable collective metabolic rates despite massive
  variation in individual activity levels. Individual ants can slow to
  near-dormancy while colony-level computation continues (Waters et al.
  2010). \emph{AI analog}: Resource allocation that maintains collective
  function under capacity constraints, graceful degradation that doesn't
  appear as degradation at the collective level.
\end{enumerate}
\end{frame}

\begin{frame}
\end{frame}

\subsection{References}\label{sec:eusocial-references}

\begin{frame}[allowframebreaks]{References}
The eusocial-intelligence literature this supplement draws on is cited
through the main bibliography (\cref{sec:references}) rather than
repeated as a second, hand-maintained list, so there is one record per
work and no opportunity for the two to disagree:

\end{frame}

\begin{frame}[allowframebreaks]{References}
\begin{itemize}
\tightlist
\item
  Foundational treatment of eusociality: \cite{wilson1971insect}.
\item
  Comprehensive ant biology: \cite{holldobler1990ants}.
\item
  Computational swarm intelligence: \cite{bonabeau1999swarm}.
\item
  Collective decision-making in bee swarms: \cite{seeley2010honeybee}.
\item
  The original stigmergy concept: \cite{grasse1959reconstruction}.
\item
  Chemical ecology and social parasitism in ants:
  \cite{lenoir2001chemical}.
\item
  Brood parasitism across insects and birds: \cite{kilner2011cuckoos}.
\end{itemize}
\end{frame}

\begin{frame}
\end{frame}

\subsection{Proofs}\label{sec:eusocial-proofs}

\subsubsection{Proof of the Redundancy-Resilience
Theorem}\label{sec:proof-redundancy-resilience}

\begin{frame}[allowframebreaks]{Proof of the Redundancy-Resilience
Theorem}
\begin{proof}
Consider a stigmergic operator $\mathcal{O}_\Sigma$ with $n$ agents, of which fraction $f$ are Byzantine (adversary-controlled).

The collective function $\mathcal{F}_c$ can be decomposed into information contributed by each agent. Let $I_i$ denote the information contribution of agent $a_i$ to the collective computation.

For the collective function to be preserved, the honest agents must contribute sufficient information:
\begin{equation}
\sum_{i \in \text{honest}} I_i \geq H(\mathcal{F}_c)
\label{eq:honest-info-sufficiency}
\end{equation}

Each honest agent contributes at most $H_{\max}$ bits. With $(1-f)n$ honest agents:
\begin{equation}
(1-f) \cdot n \cdot H_{\max} \geq H(\mathcal{F}_c)
\label{eq:honest-agent-capacity}
\end{equation}

Additionally, Byzantine consensus requires honest majority for any voting-based aggregation:
\begin{equation}
(1-f)n > 2fn \implies f < \frac{1}{3}
\label{eq:byzantine-honest-majority}
\end{equation}

Combining these constraints:
\begin{equation}
f < \min\left(\frac{1}{3}, 1 - \frac{H(\mathcal{F}_c)}{n \cdot H_{\max}}\right) = \frac{1}{3} \cdot \left(1 - \frac{H(\mathcal{F}_c)}{n \cdot H_{\max}}\right)
\label{eq:byzantine-info-combined}
\end{equation}

where the final equality holds when the information constraint is binding (typical for complex collective functions). \qed
\end{proof}
\end{frame}

\begin{frame}
\emph{This supplementary material extends the Cognitive Integrity
Framework to collective phenomena, establishing colony cognitive
security as a distinct research direction with formal foundations and
practical benchmarks.}
\end{frame}

\end{document}
